\chapter{Moteur d'exécution}
\label{chap:moteur}

\section{Principe d'exécution}

Le moteur d'exécution parcourt le plan produit par le planificateur et maintient un \textbf{contexte} : une liste de tuples partiels représentant les liaisons variable $\rightarrow$ nœud JDM découvertes jusqu'à présent.

\section{Trois cas de traitement d'une clause de relation}

\subsection{Cas 1 : Ancrage initial (Constante $\rightarrow$ Variable)}

La variable cible n'est pas encore ancrée. Le moteur interroge l'API JDM (\code{/from/} ou \code{/to/}) pour obtenir la liste des nœuds satisfaisant la clause, et crée de nouveaux tuples.

\textit{Exemple :} \texttt{(\$x r\_isa animal)} $\rightarrow$ appel à \code{/v0/relations/to/animal?types\_ids=6}. \\
Résultat : \texttt{[\{\$x: \{id: 39642, name: "chat"\}\}, \{\$x: \{id: ..., name: "chien"\}\}, ...]}

\subsection{Cas 2 : Vérification (Variable ancrée $\rightarrow$ Constante)}

La variable source est déjà ancrée. Le moteur récupère la liste valide depuis l'API, construit un \code{Set} d'IDs, et filtre le contexte en O(1).

\textit{Exemple :} \texttt{(\$x r\_has\_part queue)} après ancrage de \var{x} $\rightarrow$ appel unique + filtrage local.

\subsection{Cas 3 : Jointure (Variable $\rightarrow$ Variable)}

Les deux variables sont des variables, au moins l'une étant déjà ancrée. Le moteur itère sur les candidats de la variable ancrée, effectue un appel API par candidat (via \code{from\_by\_id} ou \code{to\_by\_id}), et conserve les couples valides.

\section{Comparaison par ID JDM}

Toutes les intersections et jointures utilisent les \textbf{identifiants numériques} des nœuds JDM, et non leurs noms textuels. Cette stratégie garantit des comparaisons exactes en O(1) avec un \code{Set<id>}, et évite les ambiguïtés dues aux homonymes ou aux variantes typographiques.

\begin{lstlisting}[caption=Comparaison par ID dans le moteur d'exécution]
// Construction d'un Set d'IDs valides (en O(n))
const validIds = new Set(apiResults.map(r => r.id));

// Filtrage local en O(1) par ID
return contexte.filter(tuple =>
    tuple[$x] && validIds.has(tuple[$x].id)
);
\end{lstlisting}

\section{Scores et preuves}

Chaque tuple porte deux métadonnées :
\begin{itemize}
    \item \code{\_\_score} : somme des poids des relations traversées ;
    \item \code{\_\_preuves} : liste des clauses qui justifient l'inclusion de ce tuple dans le résultat.
\end{itemize}

Les résultats finaux sont triés par score décroissant.

\section{Traitement des opérateurs OU}

Pour les branches OU, le moteur exécute chaque branche indépendamment et calcule l'union des résultats. Les doublons (même combinaison d'IDs) sont éliminés en conservant le tuple de score maximal.
